\section{Background and Related Work}\label{sec:background}

\subsection{How providers retire models}

OpenAI, Anthropic and Google each publish a deprecation page that lists retired and retiring model IDs. For each model, OpenAI and Anthropic give the day the deprecation was announced, the day the model stops answering (the \emph{shutdown}), and a recommended replacement. Google's page gives no announcement day and, for many models, only the earliest day on which the model may be shut down. Providers sometimes postpone a retirement; we count from the earliest announcement. A model can also have several IDs, such as a dated ID and an alias that points to it.

\subsection{Deprecation and breaking changes}

How clients react when an API they use is deprecated or changed is well studied. \citet{robbes2012deprecation} followed 577 deprecated methods and 186 deprecated classes through the Squeak and Pharo ecosystems and measured how often, how widely and how quickly client systems adapted to them, and \citet{hora2018react} measured the reactions to 118 API changes across about 3,600 Pharo systems. \citet{sawant2018deprecation} repeated the deprecation study for more than 25,000 GitHub clients of five popular Java APIs and for the JDK, and found reactions strikingly similar to those in Smalltalk. Deprecation has since been studied for Android \citep{li2018deprecated}, for Python libraries, where library maintainers handle it poorly and so make it harder for clients to move away \citep{wang2020python}, for whole Python packages \citep{zhong2024package}, and for machine-learning libraries, whose deprecated APIs developers rarely prioritise \citep{haryono2021mlapi}. LLMs themselves often suggest deprecated library APIs when completing code \citep{wang2025llmdeprecated}.

Closest to our setting are web APIs, where the provider, not the client, decides when the old version stops working. Interviews with client developers and the history of four large web APIs show how provider-driven changes impose work on client code \citep{espinha2014webapi}, and a study of 2,224 OpenAPI specifications of 1,368 RESTful APIs found several severe problems in how they deprecate \citep{yasmin2020rest}. Breaking changes in library releases are common: about one release in three in Maven broke at least one API \citep{raemaekers2017semver}, although a later differentiated replication found that only 7.9\% of clients are affected \citep{ochoa2022breaking}; in npm, about 12\% of dependent packages were hit by a breaking change, 44\% of them in minor or patch releases \citep{venturini2023broke}. Developers break APIs mostly to add features and simplify interfaces \citep{brito2018break}.

\subsection{Migration and technical lag}

Library migration, the replacement of one library by another, has been studied at scale in Java \citep{he2021migration} and Python \citep{islam2024pymig}. Projects often keep outdated dependencies, even after a security advisory \citep{kula2018update}. \emph{Technical lag} measures how far a deployed dependency is behind the ideal version \citep{zerouali2019lag}, and it has been traced across the npm dependency network \citep{decan2018lag}. Bots such as Dependabot reduce lag, but developers often tune them to send fewer updates, and some projects drop them \citep{he2023dependabot}. Pre-trained models are becoming a new kind of dependency \citep{yasmin2025deps20}, and model reuse is common in ML-enabled systems \citep{sens2025integration}. A retired LLM is a dependency with an extreme form of lag: once the shutdown day passes, the old version does not merely lag, it stops answering, so doing nothing is not an option.

\subsection{Software built on LLM APIs}

A growing body of work studies software that calls LLMs. Developers who write prompts build mental models of how the model behaves, and these models are hard to make reliable \citep{liang2025prompts}. Only 21.9\% of prompt changes in open-source repositories are documented in their commit messages \citep{tafreshipour2025prompting}, and prompts lack the testing and linting that code gets \citep{pister2024promptset}. Developers of LLM applications ask about prompts, APIs, cost and non-deterministic outputs \citep{chen2025llmdev,chen2026openaiso}, and many such questions remain unanswered \citep{alam2024challenges}. In open-source LLM applications, integration defects are common \citep{shao2025integrated}, and model issues are the most frequent kind of reported issue \citep{cai2025llmissues}; LLM-generated code also misuses library APIs \citep{zhuo2026apimisuse}.

Models behind the same API name can change behaviour over time \citep{chen2024chatgpt}, and prompts that work on one model can get worse on its successor, which calls for regression tests specific to LLM APIs \citep{ma2024prompt}. Model updates can cause regressions in generated code \citep{abbassi2025recatcher}. Yet testing practice lags behind: in ML projects, only a small share of ML-specific properties is tested often \citep{openja2025testingml}, and student teams building LLM-powered systems relied largely on exploratory testing and prompt iteration \citep{magalhaes2025untestable}. A migration to a new model is a deliberate change of exactly this kind, and it is the moment such tests would matter most. To our knowledge, apart from the study we replicate, no study has measured how open-source projects migrate between LLM versions.

\subsection{Mining commits}

Mining GitHub has well-known perils, such as forks, personal projects and inactive repositories \citep{kalliamvakou2014perils}. Labels derived from mined data can be wrong: a third of issue reports labelled as bugs are not bugs \citep{herzig2013bug}, and about 44\% of the commit messages studied in five projects could be improved \citep{tian2022commit}. We therefore validate our label on hand-labelled samples (Section~\ref{sec:method}). Bots also commit code, and can be detected from their commits \citep{dey2020bots} and their comments \citep{golzadeh2021bots}. AI coding agents are a newer source of commits: their traces suggest that between 22\% and 29\% of the 128,018 projects in one study use them \citep{robbes2026agentic}, and bot-account lookups find only a small fraction of agent commits \citep{khosravani2026agents}. This is why we keep commits with AI-tool trailers in our main label and report the result without them separately.

\subsection{The study we replicate}

\citet{kim2026retires} mined GitHub commits that migrate away from retired LLMs and reported their timing relative to each model's deprecation and shutdown, which files they change, and how they mention evaluations. It reports \KimCommits{} migration commits from \KimRepos{} repositories, of which \KimAfterShutdown\% come after the shutdown; \KimTouchTest\% of the \KimFileLevel{} migrations with file data touch a test file (Section~4.4 of that paper), and \KimMentionEval\% of commit messages mention an evaluation or regression suite (Section~4.5). It keeps commits with a migration word, drops moves into the retired model, commits from bots, commits with an AI coding-tool trailer, and commits made more than 30 days before the deprecation notice. Our replication keeps those rules as one of its labels (Section~\ref{sec:labels}) and adds a stricter one.
